\chapter{Operating Systems Cheat Sheets}

\section{Introduction}

This chapter is designed for rapid interview revision and last-minute preparation. Every section focuses on concise recall, common interview questions, key comparisons, and essential formulas.

\section{Operating System Fundamentals}

\subsection*{Key Definitions}

\begin{itemize}
\item Operating System: Resource manager and abstraction layer.
\item Kernel: Core component of the OS.
\item System Call: Interface between user programs and kernel.
\item Interrupt: Signal requiring CPU attention.
\item Trap: Software-generated interrupt.
\item DMA: Direct Memory Access.
\end{itemize}

\subsection*{Important Facts}

\begin{itemize}
\item OS manages CPU, memory, storage, and I/O devices.
\item User programs cannot directly access hardware.
\item Kernel mode has full privileges.
\item User mode is restricted.
\end{itemize}

\subsection*{Most Common Interview Questions}

\begin{enumerate}
\item What is an operating system?
\item What is a kernel?
\item What is a system call?
\item User mode vs kernel mode?
\item What is context switching?
\end{enumerate}

\subsection*{Memory Trick}

OS = Resource Manager + Abstraction Layer

\section{Computer Architecture}

\subsection*{Key Definitions}

\begin{itemize}
\item CPU
\item Register
\item Program Counter
\item Stack Pointer
\item Cache
\item MMU
\end{itemize}

\subsection*{Memory Hierarchy}

[
Registers > L1 > L2 > L3 > RAM > SSD > HDD
]

\subsection*{Interview Favorites}

\begin{itemize}
\item Fetch-Decode-Execute Cycle
\item Cache Memory
\item Interrupt Handling
\item DMA
\item NUMA
\end{itemize}

\section{Processes}

\subsection*{Key Definitions}

\begin{itemize}
\item Process = Program in execution.
\item PCB = Process Control Block.
\item Context Switch = Switching CPU execution.
\end{itemize}

\subsection*{Process States}

[
New \rightarrow Ready \rightarrow Running \rightarrow Waiting \rightarrow Terminated
]

\subsection*{Common Questions}

\begin{itemize}
\item fork vs exec
\item Zombie process
\item Orphan process
\item Context switching
\end{itemize}

\section{Threads}

\subsection*{Key Definitions}

\begin{itemize}
\item Thread = Lightweight execution unit.
\item Shares process memory.
\item Own stack and registers.
\end{itemize}

\subsection*{Benefits}

\begin{itemize}
\item Responsiveness
\item Resource sharing
\item Scalability
\item Economy
\end{itemize}

\section{CPU Scheduling}

\subsection*{Algorithms}

\begin{itemize}
\item FCFS
\item SJF
\item SRTF
\item Round Robin
\item Priority
\item Multilevel Queue
\end{itemize}

\subsection*{Metrics}

\begin{itemize}
\item Turnaround Time
\item Waiting Time
\item Response Time
\item Throughput
\item CPU Utilization
\end{itemize}

\subsection*{Formula Summary}

[
Turnaround = Completion - Arrival
]

[
Waiting = Turnaround - Burst
]

\section{Synchronization}

\subsection*{Goals}

\begin{itemize}
\item Mutual Exclusion
\item Progress
\item Bounded Waiting
\end{itemize}

\subsection*{Common Questions}

\begin{itemize}
\item Race Condition
\item Critical Section
\item Semaphore
\item Mutex
\item Monitor
\end{itemize}

\section{Critical Section Problem}

\subsection*{Requirements}

\begin{enumerate}
\item Mutual Exclusion
\item Progress
\item Bounded Waiting
\end{enumerate}

\subsection*{Classic Solutions}

\begin{itemize}
\item Peterson's Algorithm
\item Bakery Algorithm
\item Test-and-Set
\item Compare-and-Swap
\end{itemize}

\section{Semaphores}

\subsection*{Definition}

Integer synchronization variable.

\subsection*{Operations}

\begin{itemize}
\item wait()
\item signal()
\end{itemize}

\subsection*{Types}

\begin{itemize}
\item Binary Semaphore
\item Counting Semaphore
\end{itemize}

\section{Mutexes}

\subsection*{Definition}

Lock mechanism providing ownership semantics.

\subsection*{Key Difference}

Mutex owner must unlock.

\section{Deadlocks}

\subsection*{Necessary Conditions}

\begin{enumerate}
\item Mutual Exclusion
\item Hold and Wait
\item No Preemption
\item Circular Wait
\end{enumerate}

\subsection*{Strategies}

\begin{itemize}
\item Prevention
\item Avoidance
\item Detection
\item Recovery
\end{itemize}

\section{Memory Management}

\subsection*{Key Terms}

\begin{itemize}
\item Logical Address
\item Physical Address
\item Fragmentation
\item Compaction
\end{itemize}

\subsection*{Fragmentation}

\begin{itemize}
\item Internal
\item External
\end{itemize}

\section{Paging}

\subsection*{Formula}

[
PhysicalAddress=(FrameNumber\times FrameSize)+Offset
]

\subsection*{Key Terms}

\begin{itemize}
\item Page
\item Frame
\item Page Table
\item TLB
\end{itemize}

\section{Segmentation}

\subsection*{Address Format}

[
(segment, offset)
]

\subsection*{Characteristics}

\begin{itemize}
\item Logical View
\item Variable Sized
\item External Fragmentation
\end{itemize}

\section{Virtual Memory}

\subsection*{Key Concepts}

\begin{itemize}
\item Demand Paging
\item Page Fault
\item Thrashing
\item Swap Space
\item Copy-on-Write
\end{itemize}

\subsection*{Interview Favorite}

What happens during a page fault?

\section{Page Replacement}

\subsection*{Algorithms}

\begin{itemize}
\item FIFO
\item Optimal
\item LRU
\item LFU
\item Clock
\end{itemize}

\subsection*{Key Fact}

Optimal is theoretical benchmark.

\section{Memory Allocation}

\subsection*{Algorithms}

\begin{itemize}
\item First Fit
\item Best Fit
\item Worst Fit
\item Next Fit
\item Buddy System
\item Slab Allocator
\end{itemize}

\section{File Systems}

\subsection*{Key Terms}

\begin{itemize}
\item File
\item Directory
\item Inode
\item Superblock
\end{itemize}

\subsection*{Allocation Methods}

\begin{itemize}
\item Contiguous
\item Linked
\item Indexed
\end{itemize}

\section{File System Implementation}

\subsection*{On-Disk Structures}

\begin{itemize}
\item Boot Block
\item Superblock
\item Inodes
\item Data Blocks
\end{itemize}

\subsection*{Consistency}

\begin{itemize}
\item Journaling
\item fsck
\end{itemize}

\section{Disk Scheduling}

\subsection*{Algorithms}

\begin{itemize}
\item FCFS
\item SSTF
\item SCAN
\item C-SCAN
\item LOOK
\item C-LOOK
\end{itemize}

\section{I/O Systems}

\subsection*{I/O Types}

\begin{itemize}
\item Programmed I/O
\item Interrupt-Driven I/O
\item DMA
\end{itemize}

\subsection*{Models}

\begin{itemize}
\item Blocking
\item Non-Blocking
\item Asynchronous
\end{itemize}

\section{Linux Internals}

\subsection*{Must Know Topics}

\begin{itemize}
\item Process Creation
\item fork()
\item exec()
\item CFS Scheduler
\item Signals
\item systemd
\item proc
\item sysfs
\end{itemize}

\section{Modern Concurrency}

\subsection*{Key Topics}

\begin{itemize}
\item Lock-Free Programming
\item Atomic Operations
\item CAS
\item Memory Ordering
\item Thread Pools
\end{itemize}

\section{Containers}

\subsection*{Core Concepts}

\begin{itemize}
\item Namespaces
\item cgroups
\item Images
\item Containers
\end{itemize}

\section{Virtualization}

\subsection*{Key Concepts}

\begin{itemize}
\item Hypervisor
\item Type 1
\item Type 2
\item Hardware Virtualization
\end{itemize}

\section{Operating System Security}

\subsection*{Key Topics}

\begin{itemize}
\item Authentication
\item Authorization
\item Access Control
\item ASLR
\item Secure Boot
\item Sandboxing
\end{itemize}

\section{High-Concurrency Systems}

\subsection*{Must Know}

\begin{itemize}
\item Thread Pools
\item Event Loops
\item epoll
\item Reactor Pattern
\item Proactor Pattern
\item Backpressure
\end{itemize}

\section{Comparison Tables}

\subsection*{Process vs Thread}

\begin{longtable}{|p{4cm}|p{5cm}|p{5cm}|}
\hline
Aspect & Process & Thread \
\hline
Memory & Separate & Shared \
\hline
Creation Cost & High & Low \
\hline
Context Switch & Expensive & Cheaper \
\hline
Isolation & Strong & Weak \
\hline
\end{longtable}

\subsection*{Mutex vs Semaphore}

\begin{longtable}{|p{4cm}|p{5cm}|p{5cm}|}
\hline
Aspect & Mutex & Semaphore \
\hline
Ownership & Yes & No \
\hline
Counter & No & Yes \
\hline
Unlocker & Owner Only & Any Thread \
\hline
\end{longtable}

\subsection*{Paging vs Segmentation}

\begin{longtable}{|p{4cm}|p{5cm}|p{5cm}|}
\hline
Aspect & Paging & Segmentation \
\hline
Unit Size & Fixed & Variable \
\hline
Fragmentation & Internal & External \
\hline
View & Physical & Logical \
\hline
\end{longtable}

\subsection*{VM vs Container}

\begin{longtable}{|p{4cm}|p{5cm}|p{5cm}|}
\hline
Aspect & VM & Container \
\hline
Kernel & Separate & Shared \
\hline
Startup & Slow & Fast \
\hline
Isolation & Strong & Moderate \
\hline
\end{longtable}

\subsection*{FCFS vs SJF vs RR}

\begin{longtable}{|p{4cm}|p{4cm}|p{4cm}|p{4cm}|}
\hline
Aspect & FCFS & SJF & RR \
\hline
Preemptive & No & Optional & Yes \
\hline
Fairness & Medium & Poor & High \
\hline
Response Time & Poor & Good & Good \
\hline
\end{longtable}

\subsection*{Blocking vs Non-Blocking I/O}

\begin{longtable}{|p{5cm}|p{5cm}|p{5cm}|}
\hline
Aspect & Blocking & Non-Blocking \
\hline
Waits & Yes & No \
\hline
CPU Usage & Low & Higher \
\hline
Concurrency & Lower & Higher \
\hline
\end{longtable}

\subsection*{Synchronization Mechanisms}

\begin{longtable}{|p{4cm}|p{4cm}|p{4cm}|p{4cm}|}
\hline
Mechanism & Ownership & Blocking & Use Case \
\hline
Mutex & Yes & Yes & Critical Sections \
\hline
Semaphore & No & Yes & Resource Counting \
\hline
Spinlock & Yes & No & Short Waits \
\hline
Monitor & Yes & Yes & High-Level Sync \
\hline
\end{longtable}

\section{Top 100 One-Line Facts}

\begin{enumerate}
\item OS is a resource manager.
\item Kernel runs in privileged mode.
\item Threads share address space.
\item Processes provide isolation.
\item TLB caches translations.
\item Page faults trigger OS intervention.
\item Context switches are expensive.
\item Mutex provides ownership.
\item Semaphore is a counter.
\item Deadlock requires four conditions.
\item Paging eliminates external fragmentation.
\item Segmentation reflects logical program structure.
\item Virtual memory enables large address spaces.
\item LRU approximates optimal replacement.
\item Linux uses copy-on-write fork.
\item ext4 uses journaling.
\item DMA bypasses CPU copying.
\item epoll scales well.
\item Containers share kernels.
\item Hypervisors run VMs.
\item CFS schedules fairly.
\item Signals provide notifications.
\item ASLR randomizes memory layout.
\item NUMA affects memory latency.
\item Thrashing destroys performance.
\item Buddy allocators split powers of two.
\item Slab allocators cache kernel objects.
\item Working set prevents thrashing.
\item Page tables map addresses.
\item Inodes store metadata.
\item SCAN resembles an elevator.
\item SSTF minimizes seek distance.
\item Round Robin improves responsiveness.
\item Starvation affects priority scheduling.
\item Aging reduces starvation.
\item Critical sections require protection.
\item CAS supports lock-free programming.
\item Monitors encapsulate synchronization.
\item Namespaces isolate containers.
\item cgroups enforce limits.
\item Docker uses namespaces and cgroups.
\item Linux boot starts with firmware.
\item System calls cross protection boundaries.
\item Cache is faster than RAM.
\item Registers are fastest storage.
\item User mode is restricted.
\item Kernel mode is privileged.
\item File descriptors identify open files.
\item IPC enables process communication.
\item Shared memory is fastest IPC.
\item Pipes are unidirectional.
\item Sockets support networking.
\item Fork duplicates a process.
\item Exec replaces process image.
\item Zombie retains exit status.
\item Orphan is adopted by init.
\item Daemons run in background.
\item Virtualization improves isolation.
\item Containers improve density.
\item Event loops scale efficiently.
\item Thread pools reduce overhead.
\item Connection pools improve throughput.
\item Backpressure prevents overload.
\item Reactor handles readiness events.
\item Proactor handles completion events.
\item Memory-mapped files use virtual memory.
\item Swap extends memory.
\item Belady's anomaly affects FIFO.
\item Optimal replacement is theoretical.
\item Page size impacts performance.
\item Huge pages reduce TLB misses.
\item Kernel panic is fatal.
\item proc exposes kernel data.
\item sysfs exposes device information.
\item Scheduler chooses runnable tasks.
\item Locks serialize access.
\item Deadlocks stop progress.
\item Recovery breaks cycles.
\item Distributed locks extend mutexes.
\item Strong consistency returns latest writes.
\item Eventual consistency converges later.
\item Replication improves availability.
\item Consensus enables agreement.
\item Leader election selects coordinator.
\item Vector clocks track causality.
\item NTP synchronizes clocks.
\item Authentication verifies identity.
\item Authorization verifies permissions.
\item Secure boot validates startup.
\item SELinux enforces policies.
\item Cache coherence matters in multicore CPUs.
\item NUMA favors local memory access.
\item CPU scheduling affects responsiveness.
\item Throughput measures completed work.
\item Latency measures delay.
\item Availability measures uptime.
\item Reliability measures correctness.
\item Scalability measures growth capability.
\item Concurrency differs from parallelism.
\item Synchronization prevents races.
\item Protection prevents misuse.
\item Isolation improves reliability.
\item OS concepts underpin distributed systems.
\end{enumerate}

\section{Top 100 Interview Rapid Fire Questions}

\begin{enumerate}
\item What is an OS?
\item What is a kernel?
\item What is a system call?
\item What is a process?
\item What is a thread?
\item Process vs thread?
\item What is PCB?
\item What is context switching?
\item What is a race condition?
\item What is a semaphore?
\item What is a mutex?
\item What is a monitor?
\item What is deadlock?
\item Four deadlock conditions?
\item What is starvation?
\item What is aging?
\item What is paging?
\item What is segmentation?
\item What is TLB?
\item What is a page fault?
\item What is virtual memory?
\item What is thrashing?
\item What is swap?
\item What is copy-on-write?
\item What is FIFO?
\item What is LRU?
\item What is Optimal replacement?
\item What is fragmentation?
\item Internal fragmentation?
\item External fragmentation?
\item What is buddy allocation?
\item What is slab allocation?
\item What is an inode?
\item What is a superblock?
\item What is journaling?
\item What is ext4?
\item What is DMA?
\item What is interrupt-driven I/O?
\item What is blocking I/O?
\item What is non-blocking I/O?
\item What is asynchronous I/O?
\item What is select?
\item What is poll?
\item What is epoll?
\item What is Linux CFS?
\item What is fork()?
\item What is exec()?
\item What is a zombie process?
\item What is an orphan process?
\item What is a signal?
\item What is proc?
\item What is sysfs?
\item What is a namespace?
\item What is a cgroup?
\item What is Docker?
\item What is Kubernetes?
\item What is a hypervisor?
\item Type 1 hypervisor?
\item Type 2 hypervisor?
\item What is virtualization?
\item What is a VM?
\item What is a container?
\item What is ASLR?
\item What is secure boot?
\item What is SELinux?
\item What is NUMA?
\item What is cache coherence?
\item What is a spinlock?
\item What is CAS?
\item What is memory ordering?
\item What is a thread pool?
\item What is a connection pool?
\item What is backpressure?
\item What is a reactor?
\item What is a proactor?
\item What is throughput?
\item What is latency?
\item What is scalability?
\item What is availability?
\item What is consistency?
\item What is consensus?
\item What is Raft?
\item What is Paxos?
\item What is a distributed lock?
\item What is leader election?
\item What is NTP?
\item What is a logical clock?
\item What is a vector clock?
\item What is CAP theorem?
\item Strong consistency?
\item Eventual consistency?
\item What is replication?
\item What is fault tolerance?
\item What is 2PC?
\item What is IPC?
\item Shared memory?
\item Pipe?
\item Socket?
\item Why does thread-per-request not scale?
\item How does Nginx scale?
\end{enumerate}

\section{Final Interview Day Checklist}

\begin{itemize}
\item Revise Process vs Thread.
\item Revise Mutex vs Semaphore.
\item Memorize Deadlock Conditions.
\item Review Page Fault Handling.
\item Review TLB and Paging.
\item Review Virtual Memory.
\item Review Scheduling Algorithms.
\item Review Synchronization Problems.
\item Review Linux Process Lifecycle.
\item Review Containers and Virtualization.
\item Review epoll and Event Loops.
\item Practice concise answers.
\item Focus on trade-offs.
\item Explain concepts visually.
\item Connect OS to system design.
\end{itemize}

\section{Complete OS Revision Map}

\begin{enumerate}
\item Fundamentals
\item Computer Architecture
\item Processes
\item Threads
\item Scheduling
\item Synchronization
\item Deadlocks
\item Memory Management
\item Paging
\item Segmentation
\item Virtual Memory
\item Page Replacement
\item Allocation
\item File Systems
\item Storage
\item I/O
\item Linux
\item Security
\item Virtualization
\item Containers
\item Concurrency
\item Distributed Systems Connections
\item High-Concurrency Systems
\item System Design Applications
\end{enumerate}

\section{Key Takeaways}

\begin{keytakeaways}
Master the core concepts before memorizing advanced details. Focus heavily on processes, threads, synchronization, deadlocks, virtual memory, paging, Linux internals, containers, and high-concurrency architectures. Understand trade-offs rather than definitions alone. Practice explaining concepts clearly, drawing diagrams, solving numerical problems, and connecting operating-system fundamentals to real-world backend systems and large-scale system design.
\end{keytakeaways}
